\section{Method}\label{sec:method}

\subsection{Platforms and Sessions}\label{sec:platforms}
\begin{table*}[t]
\caption{\capTabPlatforms}\label{tab:platforms}
\centering\small\gentable{tab_platforms}\gennote{tab_platforms}
\end{table*}

A platform is one of our three hardware targets. These are the TPU v5e chip, the NVIDIA T4 GPU, and the Google Colab
CPU runtime. A session is one Colab runtime or one Kubernetes Job in which a set of runs took place.
In the September sessions, the CPU runtime had two logical CPUs. % src: results/repro/2026-09-15_cpu-colab/environment_cpu-colab-repro.json -> hardware.logical_cpus; recomputed: SEP.cpu.0915.logical_cpus
Table~\ref{tab:platforms} lists every session with a baseline run.

The TPU runs used a Google Kubernetes Engine cluster of the course, in the us-west4 region. Each Kubernetes Job
requested the whole slice of eight chips. Some Jobs ran several experiments.
The committed Jobs started from a \texttt{python:3.12-slim} image and installed \texttt{jax[tpu]==0.10.2}. % src: configs/af_spike_job.yaml ("image: python:3.12-slim"); recomputed: V2.jax_pin
They pinned no other package. They cloned AlphaFold2 from its default branch and did not record the commit. The TPU
runs did not record their dates.
Most TPU results are in the first commit of our benchmark repository, dated 8 August 2026, 17:23 UTC. % src: recomputed: V2.first_commit
We committed the results of the matched grid (Section~\ref{sec:res-pmap}) on the same day, at 18:05 UTC, and those of the external ensemble on 11 August 2026. % src: recomputed: V2.grid_commit, V2.ensemble_commit
We cannot identify the script of the baseline run on the TPU. Its result file has all but one of the fields that the
notebook copy writes. It has none of the extra fields that the committed script writes. % src: results/result_tpu-v5e-podslice.json -> keys; results/result_cpu-colab.json -> keys; recomputed: B.tpu.has_host_processor_field

The baseline runs on the T4 and the CPU runtime ran in two Colab runtimes on 8 August 2026, at 04:50 and 05:19 UTC. % src: recomputed: AUG.gpu.first_epoch_utc, AUG.cpu.first_epoch_utc
They used a copy of the benchmark script inside the notebooks.
The notebooks in commit 2755250 hold this copy and the output of the runs. This copy differs from the scripts in \texttt{src/}. % src: recomputed: V2.first_commit; git show 2755250:notebooks/alphafold_cpu_benchmark.ipynb, cell 8 (script_content, output "Wrote 7449 bytes")
The first calls of these two runs ran with the profiler. Table~\ref{tab:platforms} does not show them. These two sessions did not
record their JAX version or their host CPU. In September, we ran the committed script on the T4 and the CPU runtime.
The CPU runtime ran in three sessions, on 15, 16, and 17 September 2026. The T4 ran in one session, on 15 September 2026. % src: recomputed: V2.sep_dates
Each session ran the script twice, once with the profiler and once without it. For the September sessions, tables and
text report the runs without the profiler. Section~\ref{sec:res-first} also uses the runs with the profiler. Each
September session recorded its commit, its package versions, and its host CPU.
These sessions ran commits 86fca03 and 9e1a07e, which have the same benchmark scripts. % src: recomputed: V2.sep_src_identical; results/repro/2026-09-15_gpu-t4/environment_gpu-t4-repro.json -> repo_commit; results/repro/2026-09-16_cpu-colab/environment_cpu-colab-repro.json -> repo_commit
All four sessions ran JAX 0.10.2 and AlphaFold2 at commit c77e5d2~\cite{alphafoldrepo}. % src: recomputed: SEP.cpu.0915.jax, SEP.cpu.0916.jax, SEP.cpu.0917.jax, SEP.gpu.0915.jax, SEP.cpu.0915.alphafold_commit
We did not repeat the TPU runs. Our access to the course cluster ended with the course.
% The plan of the repeat is in commit 1651c42 (18 September 2026).
We prepared a repeat on Cloud TPU in the folder \texttt{cloud/}. We had no cloud resources when we prepared it, and
we did not run it later.

\subsection{Workload and Inputs}\label{sec:workload}
Every run except the model comparison used \texttt{model\_3} of AlphaFold2. This model needs no templates. Templates are known
structures of related proteins that other AlphaFold2 models take as input. The script built the model with random
weights. Haiku creates these weights in the parameter initialization call. The script never loaded the trained
checkpoint. The MSA held only the query sequence, and no run searched a sequence database. The featurization pads the
MSA features to a fixed number of rows.
With one sequence, \texttt{model\_3} still runs on 512 MSA rows and 5120 extra MSA rows~\cite{alphafoldrepo}. % src: AlphaFold2 at commit c77e5d2, alphafold/model/config.py:233 ("'max_msa_clusters': 512"), config.py:73 ("'data.common.max_extra_msa': 5120"), alphafold/model/tf/data_transforms.py:389-425 (make_fixed_size); paper-v3/SOURCE_NOTES.md section 3
Only the first MSA row holds the query, and the other rows are padding.
Every run used no recycling, one ensemble member, float32 parameters, and random seed 0, unless an experiment varied one of them. % src: src/spike_tpu_forward_pass.py ("random_seed=0")
The baseline configuration is the single-query script on the fixed sequence with \texttt{model\_3}, no recycling, and
float32 parameters. It times \texttt{predict}. AlphaFold2's source ran unmodified. Newer JAX releases changed the
arguments of one JAX function that AlphaFold2 calls. A small wrapper outside AlphaFold2 renamed these arguments.

The experiments used two input families. The fixed sequence has 118 residues. % src: recomputed: INPUT.toy_length
The repeated-alphabet sequence repeats the 20 amino-acid letters to the target length. % src: recomputed: V2.alphabet_length, INPUT.per_experiment ("alphabet ACDEFGHIKLMNPQRSTVWY repeated to the length")
In runs with several queries, each query starts one letter later in the alphabet. Table~\ref{tab:experiments} gives
the family of each experiment.

\subsection{Timing and Memory}\label{sec:timing}
The scripts time each call with \texttt{time.time} on the host. They wait for the outputs of the first call and the
steady-state calls with \texttt{jax.block\_until\_ready}. Parameter initialization is the \texttt{init\_params} call,
which traces the model and creates the random weights. The first call is the earliest call of \texttt{predict} or
\texttt{apply} at a new input shape. It traces, compiles, and executes the model. A steady-state call is a later call at
the same input shape, which reuses the compiled program. The single-query and automatic-partitioning scripts time
parameter initialization, the first call, and the steady-state calls. The \texttt{vmap}, \texttt{pmap}, and ensemble
scripts time the first call and one steady-state call. The August Colab runs and all TPU runs timed one steady-state
call, and the September runs without the profiler timed five. % src: src/spike_tpu_forward_pass.py ("--num_steady_state_runs=5")

The scripts time two different regions. The single-query script and the automatic-partitioning script time
\texttt{predict}, which includes the confidence metrics on the host. The \texttt{vmap}, \texttt{pmap}, and ensemble
scripts time \texttt{apply} only (Table~\ref{tab:experiments}). The single-query script and the \texttt{vmap} script
run the first call inside \texttt{jax.profiler.trace}, a context that records a profile of the call. The September runs
without the profiler are the exception. The profiler teardown is the time this context takes to close after the first
call returns. It lies inside the timer of the first call.

Throughput is the number of queries completed per second in steady-state calls, in queries/s. The speedup is the
throughput on $n$ chips divided by the throughput on one chip, with the same input family and timed region. The parallel
efficiency is the speedup divided by $n$. With multi-query \texttt{pmap}, each chip runs one query, so the speedup
measures throughput and not the time of one query. On one chip, multi-query \texttt{pmap} runs one query. Every
\texttt{pmap} run had all eight chips visible.
% Every script in src/ records stats.get("bytes_in_use") of device.memory_stats() for each
% device of jax.devices(). The definition of MB is in the note of Table III (tools/make_tables.py).
HBM in use is the value of \texttt{bytes\_in\_use} that JAX reports for a chip after the first steady-state call. It
is the device memory allocated at that moment. Peak memory is the largest HBM in use over the life of the process.

\subsection{Experiments}\label{sec:experiments}
\begin{table*}[t]
\caption{\capTabExperiments}\label{tab:experiments}
\centering\small\setlength{\tabcolsep}{4pt}\gentable{tab_experiments}\gennote{tab_experiments}
\end{table*}

Table~\ref{tab:experiments} lists the baseline configuration and the 12 TPU experiments, with script, input, timed region, and number of runs. % src: recomputed: COUNT.tpu_experiments_with_raw_data
In this table, script X is the file \texttt{src/spike\_X\_forward\_pass.py}, with \texttt{tpu} for single,
\texttt{batch} for vmap, \texttt{meshshard} for partitioning, and \texttt{ensemble\_shard} for ensemble.
Experiments in one Job ran one after another in the same container, in a fixed order. Most settings ran once. The
repeated-runs experiment ran the baseline configuration three times on the TPU. Its coefficient of variation was 0.12\% for the steady-state call, 0.82\% for parameter initialization, and 1.00\% for the first call. % src: results/sweep/repeated_runs.json -> runs; recomputed: RR.steady_state_seconds.cv_pct, RR.init_params_seconds.cv_pct, RR.first_predict_seconds.cv_pct
We ran the baseline configuration 11 times on the TPU. Ten of these runs started with an empty compilation cache. % src: recomputed: TPU.single_query_steady_n, V2.tpu_first_n
We also tried to make only two or four chips visible to the single-query script.
The file \texttt{results/sweep/chip\_visibility.md} records that these runs stopped with an error. We have no log of them. % src: results/sweep/chip_visibility.md ("Attempted: 2 and 4 chip subsets")

\subsection{Records}\label{sec:records}
Most TPU results were committed together with their scripts. The files in \texttt{results/sweep/} are summaries of the
Job output. They are not the files that the scripts wrote. Most of their field names differ from those that the
scripts write. Every memory value in them is a whole number of megabytes. % src: recomputed: CV.all_bytes_multiple_of_1e6
The logs of the Jobs are not in the repository, and we have not found them. The first author wrote these summaries
from the output of the Jobs with the help of Claude (Anthropic). We have no record of how each value was transferred. We
report the values in these summaries as they are. We checked the summaries against each other. In the \texttt{vmap}
and \texttt{pmap} summaries, each throughput equals the number of queries divided by the steady-state call, within the
rounding of the steady-state call.
The 11 single-query runs of the baseline configuration come from seven files and took 0.469~s to 0.471~s. % src: recomputed: TPU.single_query_steady_n, TPU.single_query_steady_min_max; paper-v3/SOURCE_NOTES.md ("seven files")
These checks cannot replace the output of the Jobs.
